% A0 landscape poster with beamerposter: an ordinary Beamer document with one huge page, so
% blocks, columns and TikZ work as in slides. (The `poster` template uses tikzposter instead.)
% For A1 or portrait, change size= / orientation= below and adjust the column widths.
\documentclass[final]{beamer}
\usepackage[size=a0,orientation=landscape,scale=2]{beamerposter}
\usepackage[T1]{fontenc}
\usepackage{lmodern}
\usepackage{amsmath}
\usepackage{booktabs}
\usepackage{tikz}
\usetikzlibrary{arrows.meta,positioning}
\usepackage{pgfplots}
\pgfplotsset{compat=1.18}

% --- Look: one accent colour, rounded blocks, no navigation symbols. ---
\usetheme{default}
\setbeamertemplate{navigation symbols}{}
\renewcommand{\familydefault}{\sfdefault}
\definecolor{accent}{HTML}{1F5FA8}
\setbeamercolor{block title}{bg=accent, fg=white}
\setbeamercolor{block body}{bg=accent!8, fg=black}
\setbeamertemplate{blocks}[rounded][shadow=false]
\setbeamerfont{block title}{size=\LARGE, series=\bfseries}
\setbeamerfont{block body}{size=\large}
\setbeamersize{text margin left=1.5cm, text margin right=1.5cm}
\addtobeamertemplate{block end}{}{\vspace{1.2cm}} % air between blocks
\setbeamertemplate{itemize items}[circle]
\setbeamercolor{itemize item}{fg=accent}
\setlength{\parskip}{0.6ex}

% --- Logo placeholders drawn with TikZ: replace each by \includegraphics[height=5cm]{logo.pdf}. ---
\newcommand{\logoleft}{\tikz[baseline=(current bounding box.center)]{
  \draw[white, line width=2mm] (0,0) rectangle (7,7); \fill[white] (0.7,0.7) rectangle (6.3,6.3);
  \node[accent, font=\LARGE\bfseries] at (3.5,3.5) {Uni};}}
\newcommand{\logoright}{\tikz[baseline=(current bounding box.center)]{\fill[white] (3.5,3.5) circle (3.5);
  \fill[accent] (3.5,3.5) circle (2.9); \node[white, font=\LARGE\bfseries] at (3.5,3.5) {Lab};}}

% --- A QR-code-sized placeholder (argument: size of one module in cm; 9 modules wide). Swap for \includegraphics{qr.pdf} or the `qrcode` package. ---
\newcommand{\qrplaceholder}[1]{\begin{tikzpicture}[x=#1cm, y=#1cm]
  \draw[line width=0.8mm] (0,0) rectangle (9,9);
  \foreach \x/\y in {0.7/6.3, 6.3/6.3, 0.7/0.7}{
    \draw[line width=1.2mm] (\x,\y) rectangle ++(2,2); \fill (\x+0.6,\y+0.6) rectangle ++(0.8,0.8);}
  \foreach \x/\y in {4/7, 4.6/7.8, 5.3/6.6, 4/5, 5/4.4, 6/5.4, 7/4, 4.2/3, 5.4/2.2, 6.6/2.8, 7.4/1.6, 4/1.2, 3.2/4.2}{
    \fill (\x,\y) rectangle ++(0.7,0.7);}
\end{tikzpicture}}

\title{Poster Title: One Clear Message in at Most Two Lines}
% \texorpdfstring keeps the superscripts out of the PDF's author property (otherwise hyperref warns).
\newcommand{\aff}[1]{\texorpdfstring{$^{#1}$}{ }}
\author{First Author\aff{1}, Second Author\aff{2}, Third Author\aff{1}}
\institute{$^1$Example University, Department of Science \quad $^2$Example Institute\\
  \texttt{first.author@example.edu}}

% The header: title and authors in the middle, a logo at each end.
\setbeamertemplate{headline}{%
  \begin{beamercolorbox}[wd=\paperwidth, ht=17cm, dp=0pt]{title}
    \vbox to 17cm{\vfill
      \hbox to \paperwidth{\hspace{2cm}\logoleft\hfill
        \parbox[c]{0.6\paperwidth}{\centering
          {\fontsize{80}{88}\selectfont\bfseries \inserttitle\par}\vspace{0.8cm}
          {\huge \insertauthor\par}\vspace{0.3cm}
          {\Large \insertinstitute\par}}%
        \hfill\logoright\hspace{2cm}}%
    \vfill}
  \end{beamercolorbox}%
}
\setbeamercolor{title}{bg=accent, fg=white}

\begin{document}
\begin{frame}[t]
\vspace{0.5cm}
\begin{columns}[t,onlytextwidth]

  % --- Column 1: why and what ---
  \begin{column}{0.24\textwidth}
    \begin{block}{Motivation}
      Say why anyone should stop at your poster, in three short lines:
      \begin{itemize}
        \item The problem in one sentence.
        \item Why existing approaches fall short.
        \item What you did instead.
      \end{itemize}
    \end{block}
    \begin{block}{Research question}
      One question, readable from two metres away: \emph{Does the new method beat the
      baseline without extra cost?}
    \end{block}
    \begin{block}{Data}
      Describe the data in two lines: where it comes from, how big it is, what was excluded.
    \end{block}
  \end{column}

  % --- Column 2: method ---
  \begin{column}{0.24\textwidth}
    \begin{block}{Method}
      A diagram beats a paragraph; a formula should be set big enough to read from far away:
      \[ \Large \hat{\theta} = \operatorname*{arg\,min}_{\theta}\; \frac{1}{n}\sum_{i=1}^{n} \ell\bigl(f_\theta(x_i), y_i\bigr) \]
      \centering
      \begin{tikzpicture}[node distance=3cm, box/.style={draw=accent, fill=white, line width=1mm, rounded corners, minimum width=16cm, minimum height=4.5cm}]
        \node[box] (a) {Data};
        \node[box, below=of a] (b) {Model};
        \node[box, below=of b] (c) {Evaluation};
        \draw[-{Stealth[length=1cm]}, line width=1mm] (a) -- (b);
        \draw[-{Stealth[length=1cm]}, line width=1mm] (b) -- (c);
      \end{tikzpicture}
    \end{block}
  \end{column}

  % --- Column 3: results ---
  \begin{column}{0.24\textwidth}
    \begin{block}{Results}
      \centering
      \begin{tikzpicture}
        \begin{axis}[scale only axis, width=0.72\linewidth, height=32cm, xlabel={Training size (k)}, ylabel={Accuracy (\%)},
            legend pos=south east, grid=major, grid style={gray!30}, line width=1mm, mark size=4pt,
            xtick={0,4,8,12,16}, ytick={60,70,80,90}, ymin=55, ymax=92, xmin=0,
            tick label style={font=\Large}, label style={font=\Large}, legend style={font=\Large}]
          \addplot[accent, mark=*] coordinates {(1,62) (2,70) (4,77) (8,82) (16,87)};
          \addplot[gray, mark=square*, dashed] coordinates {(1,58) (2,64) (4,71) (8,76) (16,81)};
          \legend{Ours, Baseline}
        \end{axis}
      \end{tikzpicture}

      \bigskip
      \begin{tabular}{@{}lrr@{}}
        \toprule
        Method & Acc.\ (\%) & Time (s) \\
        \midrule
        Baseline & 81.2 & 12 \\
        Ours & \textbf{86.9} & 15 \\
        \bottomrule
      \end{tabular}
    \end{block}
  \end{column}

  % --- Column 4: take-home ---
  \begin{column}{0.24\textwidth}
    \begin{block}{Conclusions}
      \begin{itemize}
        \item The main finding, with its number: +5.7 points.
        \item The cost, stated honestly: +3~s.
        \item What to do next.
      \end{itemize}
    \end{block}
    \begin{block}{References}
      [1] L.~Lamport. \emph{\LaTeX: A Document Preparation System}. Addison-Wesley, 1994.\par
      [2] D.~E.~Knuth. \emph{The \TeX book}. Addison-Wesley, 1984.
    \end{block}
    \begin{block}{Paper, code and contact}
      \centering
      \qrplaceholder{1.8}\par
      \smallskip
      example.org/our-project
    \end{block}
  \end{column}

\end{columns}
\end{frame}
\end{document}
